% ===================================================================
%  नकसा नं. ९ — पहाड़। जंगल।
%
%  Traduction, faite en 2026, en bhojpouri d'aujourd'hui, sur l'ido
%  de texto/io. Regles de la colonne : voir 10-nakasa-01.tex. DEUX
%  scenes, deux numerotations qui repartent de 1 (t09-c1-*,
%  t09-c2-*). 47 blocs, 117 renvois, trente-quatre paires a
%  retourner. Le bloc t09-c1-05-1 tient DEUX alineas sous une seule
%  cle : ne pas en ouvrir une seconde.
%
%  LA PLAINE NOMME LA MONTAGNE AVEC DU SANSCRIT EMPRUNTE, EXACTEMENT
%  COMME ELLE NOMMAIT LA LUMIERE AVEC DE L'ANGLAIS. Le pays bhojpouri
%  est l'un des endroits les plus plats de la terre : de Chapra a
%  Buxar, sur deux cents kilometres, il n'y a ni rocher, ni ravin, ni
%  torrent, ni glacier. La langue a pourtant tous les mots — हिमनद le
%  glacier, ज्वालामुखी le volcan, घाटी la vallee, खाई le ravin,
%  हिमस्खलन l'avalanche, चट्टान le rocher — et ils sont tous des
%  tatsama sanskrits, pris tels quels dans une langue savante. Le
%  tableau 6 avait montre la meme mecanique sur la lumiere, avec
%  l'anglais et le persan pour preteurs. ICI LE PRETEUR EST PLUS
%  VIEUX, MAIS LA LOI EST LA MEME : une langue nomme par son fonds ce
%  qu'elle a, et par emprunt ce qu'elle n'a pas. Et la montagne n'est
%  pas inconnue — l'Himalaya est a deux cents kilometres au nord et se
%  voit de la plaine par les matins clairs de janvier. Elle n'est pas
%  ignoree : elle n'est pas habitee.
%
%  LE MONTREUR D'OURS (41, 44) A COUTE DEUX REFUS, ET LE SECOND EST
%  LE PLUS DIFFICILE DE LA COLONNE. L'ido ecrit « cigano », que le
%  francais de 1926 rend par bohemien. Le premier refus est facile :
%  « जिप्सी » est une insulte en 2026 et ne s'ecrit pas ; on ecrit
%  « रोमा ». Le second l'est beaucoup moins. Le montreur d'ours EXISTE
%  dans la plaine bhojpourie, il s'appelle मदारी ou कलंदर, il mene sa
%  bete a la chaine exactement comme celui de la planche — et la
%  langue des Roms descend d'un indo-aryen central cousin de celui
%  d'ou vient le bhojpouri, de sorte que les deux hommes exercent le
%  meme metier dans deux langues sœurs. LA TENTATION D'ECRIRE मदारी
%  EST DONC MAXIMALE, ET C'EST PRECISEMENT POUR CELA QU'IL FAUT S'EN
%  ABSTENIR : un Rom des Pyrenees n'est pas un मदारी de Chapra, meme
%  si l'un et l'autre descendent de la meme migration. On modernise la
%  langue, jamais les choses. (Le metier a cesse des deux cotes :
%  l'ours danseur est interdit en Inde depuis 1972, et le dernier a
%  ete rendu en 2009.)
%
%  Enfin le chamois (54) n'a pas de nom : la plaine n'a pas
%  d'antilope de rocher, et l'Himalaya en a d'autres — le घोरल, le
%  सीरो — qui ne sont pas celle-la. On ecrit « चामोइस के सिकारी »,
%  comme on a ecrit « बड़का हँसुआ » pour la faux au tableau 7.
% ===================================================================

%%K t09-apar-1 apar
\VUsaut{4.0mm}
\VUtitre{15.0pt}{60}{नकसा नं. ९}
\VUsaut{3.0mm}

%%K t09-tit-1 sub
\VUcentre{12.0pt}{0}{\textit{पहिलका दृस्य।}}
\VUsaut{3.0mm}

%%K t09-tit-2 sub
\VUpk{11.4pt}{40}{पहाड़। --- पानी के सहर।}
\VUsaut{3.0mm}

%%K t09-c1-01-1 p
१. --- आठ दिन से हम प्रात्स में बानी। हम ऊ पूरा अचरज ना बता सकीं जवन
हमरा भइल जब हम पिरेनी के ओह अद्भुत \VUgras{पहाड़ के कतार}
\textsuperscript{(1)} के सामने रहनी, जेकर \VUgras{कगार}
\textsuperscript{(2)} नीला आसमान पर बड़हन झालर नियर लउकत बा।

%%K t09-c1-02-1 p
२. --- हमरा सोचो में ना रहे कि पिरेनी के \VUgras{चोटी}
\textsuperscript{(3)} एतना \VUgras{ऊँचाई} तक पहुँचेला, आ हम ओह सुंदर
\VUgras{हिमनद} \textsuperscript{(5)} आ बरफ से ढँकल \VUgras{सिखर}
\textsuperscript{(4)} के सामने दंग रह गइनी, जवना पर एतना डरावन
\VUgras{हिमस्खलन} लुढ़केला।

%%K t09-tit-3 sub
\VUsaut{3.0mm}
\VUpk{10.2pt}{40}{पहाड़ के नजारा।}
\VUsaut{3.0mm}

%%K t09-c1-03-1 p
३. --- पहुँचला के अगिले दिन हम ओह बुझल \VUgras{ज्वालामुखी}
\textsuperscript{(6)} लगे गइनी जवन प्रात्स के लगे बा।
\VUgras{मुहाना} के सूखल आ नंगा किनारा पर ओह \VUgras{लावा} के बहला के
निसान अबहीं ले साफ लउकत बा, जवन कई सदी पहिले \VUgras{पहाड़ के ढलान}
\textsuperscript{(7)} पर बहल रहे। हम एगो ढलान पर ओह लावा के कुछ टुकड़ा
खोज लेनी।

%%K t09-c1-04-1 p
४. --- प्रात्स सरहद पर बसल बा, आ जवन \VUgras{किला} \textsuperscript{(8)}
ओह \VUgras{दर्रा} \textsuperscript{(27)} के रच्छा करेला जवन फ्रांस के
स्पेन से अलग करेला, ऊ राइन के किनारे के ओह पुरनका किलन के इयाद दिआवेला
जे आपन चट्टान के चोटी से \VUgras{सिकार} ताकत लागेला।

%%K t09-c1-05-1 p
५. --- एगो बड़हन \VUgras{उकाब} \textsuperscript{(9)}, जेकर घोंसला
\VUgras{किला} \textsuperscript{(8)} से दूर नइखे, अकसर हमार धेयान खींचेला।

ई \VUgras{सिकारी चिरई}, जेकर \VUgras{चोंच} मजबूत बा, अकसर आपन बच्चा
खातिर मेमना भा पठरू ले आवेला, जेकरा के ऊ आपन \VUgras{पंजा} में थामले
रहेला। ओकर \VUgras{पाँख} एतना बड़ बा कि ऊ आपन भारी बोझ के साथे देर तक
हवा में तिर सकेला।

%%K t09-tit-4 sub
\VUsaut{3.0mm}
\VUpk{10.2pt}{40}{घुमक्कड़।}
\VUsaut{3.0mm}

%%K t09-c1-06-1 p
६. --- ओहिजा सबसे नीक घुमाई « काउ » के झरना \textsuperscript{(10)} तक के
जातरा ह। \VUgras{पानी के झरना} भँवर बनावत गिरेला आ ओकरा बाद सहर के
पच्छिम में बसल \VUgras{सनोबर के जंगल} \textsuperscript{(11)} में सुंदर
\VUgras{नाला} बन के गायब हो जाला। हमनी के एह जंगल से होके \VUgras{मौसम
के वेधशाला} \textsuperscript{(12)} तक चढ़नी जा, आ \VUgras{मचान} के ऊपर से
इलाका के पूरा \VUgras{नजारा} देखनी जा। \VUgras{दृस्य} अद्भुत बा, आ
\VUgras{प्रकृति} जइसे एह देस के आपन सब खजाना लुटा देत बाड़ी।

%%K t09-c1-07-1 p
७. --- उतरे खातिर हमनी के \VUgras{रस्सा-रेल} \textsuperscript{(13)}
चलवनी जा आ एगो छोट \VUgras{इस्टेसन} पर रुकनी जा, जवन एगो पुरनका
\VUgras{किला} के \VUgras{खंडहर} \textsuperscript{(14)} लगे बा, जहाँ कबो
प्रात्स के सामंत रहत रहलें।

%%K t09-c1-08-1 p
८. --- बेचारा किला ! ऊ का सोचत होई जब नीचे ओह सजल-धजल \VUgras{पानी के
सहर} \textsuperscript{(15)} के देखत होई जवन अबहीं चलन के \VUgras{गरम पानी के
इस्टेसन} बन गइल बा ?

%%K t09-c1-08-2 p
\VUgras{जलवायु} एतना सुहावन बा, \VUgras{खनिज पानी} एतना असरदार बा, कि
\VUgras{सैनेटोरियम} \textsuperscript{(17)} में होखे वाला \VUgras{इलाज}
सचहूँ सुख के बात बा।

%%K t09-tit-5 sub
\VUsaut{3.0mm}
\VUpk{10.2pt}{40}{सुहावन गरम पानी के सहर।}
\VUsaut{3.0mm}

%%K t09-c1-09-1 p
९. --- ऊपर से, रहे के सुहावन बनावे खातिर सब कुछ कइल गइल बा। रेल खातिर
एगो बड़का \VUgras{पुलिया} \textsuperscript{(16)} बनावल गइल ; \VUgras{गरम
पानी के भवन} के \VUgras{बगइचा} \textsuperscript{(18)}, जहाँ
\VUgras{संगीत के जलसा} होला, बहुते सुंदर ढंग से सजल बा। कई गो सुंदर
« \VUgras{कोठी} » \textsuperscript{(19)} जुआघर \textsuperscript{(21)}
के चारो ओर बनल बा, आ नीक से राखल \VUgras{पक्का सड़क}
\textsuperscript{(22)} आवे-जाए के बहुते आसान बना देले बा।

%%K t09-c1-10-1 p
१०. --- एगो सादा \VUgras{ढाल} \textsuperscript{(23)} \VUgras{सड़क}
\textsuperscript{(20)} के ओह \VUgras{घाटी} \textsuperscript{(25)} से अलग
करेला जेकर तल में एगो सुंदर छोट \VUgras{झील} \textsuperscript{(26)} बा,
जवन एगो साफ \VUgras{नाला} \textsuperscript{(24)} के पानी देला।

%%K t09-c1-11-1 p
११. --- घाटी के दोसरा ओर \VUgras{लोहा के खान} \textsuperscript{(28)}
खोदल जात बा। हम कई बेर \VUgras{खनिहार} लोग के \VUgras{कुइयाँ} में उतरत
देखे गइनी, आ हम ओह \VUgras{सुरंग} में भी गइनी जहाँ ऊ लोग आपन दिन
बितावेला, ओह \VUgras{कच्चा धातु} के निकालत जवना में \VUgras{धातु}
रहेला।

%%K t09-c1-12-1 p
१२. --- खान लगे एगो बड़का \VUgras{ढलाईखाना} बनावल गइल बा ताकि ओहसे
निकलल संपत्ति के तुरते काम में लावल जाव। \VUgras{बड़का भट्ठी}
\textsuperscript{(29)}, जवन इहाँ भरपूर मिले वाला \VUgras{पथरकोयला} से
गरम होला, जल्दी आ सस्ता में \VUgras{ढालल लोहा} देवेला।

%%K t09-c1-13-1 p
१३. --- खान से दूर नइखे कि एगो \VUgras{पाथर के खान} \textsuperscript{(30)}
खोदल जात बा। बुढ़उ \VUgras{पाथर तोड़े वाला} आपन \VUgras{गैंता} से ना
\VUgras{ग्रेनाइट} काटेलें, ना \VUgras{पोरफरी}, ना \VUgras{बलुआ पाथर} —
खाली घर बनावे खातिर सादा \VUgras{गढ़ल पाथर}।

%%K t09-c1-14-1 p
१४. --- एह देस के सबसे सुंदर सिंगार में से एगो \VUgras{सनोबर}
\textsuperscript{(31)} ह। हर जगह — \VUgras{खाई} \textsuperscript{(32)} के
तल में, \VUgras{पहाड़ी धार} \textsuperscript{(33)} के किनारे,
\VUgras{गहिर खाई} \textsuperscript{(34)} भा खड़ी ढाल के लगे — ओकर पातर
पत्ता हरियर छाँह बिखेरत बा।

%%K t09-tit-6 sub
\VUsaut{3.0mm}
\VUpk{10.2pt}{40}{पैदल घुमाई।}
\VUsaut{3.0mm}

%%K t09-c1-15-1 p
१५. --- हम आपन छुट्टी के फायदा उठा के \VUgras{लमहर} घुमाई करत बानी।
जवन \VUgras{रस्ता} \textsuperscript{(35)} पहाड़ पर घुमावदार चलेला ओहसे
दूरी के परवाह कइले बिना कई गो \VUgras{किलोमीटर} आ अकसर कई गो
\VUgras{चार किलोमीटर} नाप लिहल जाला।

%%K t09-c1-15-2 p
\VUgras{मील के पाथर} \textsuperscript{(36)} आ \VUgras{दिसा-खंभा}
\textsuperscript{(37)} रस्ता के निसान देवेला, जवना पर अकसर एगो जवान
\VUgras{पहाड़ी} \textsuperscript{(38)} मिलेला, बगल में \VUgras{झोरी}
\textsuperscript{(40)} लटकवले, \VUgras{मारमोट} \textsuperscript{(39)} लेके
चलत, भा कवनो \VUgras{रोमा} \textsuperscript{(41)}, जेकर \VUgras{फर के
कोट} \textsuperscript{(42)} ओकर फाटल पुरनका \VUgras{टोपा}
\textsuperscript{(43)} से अजीब मेल खात बा। ऊ \VUgras{जंजीर} से एगो सधावल
\VUgras{भालू} \textsuperscript{(44)} ले चलत बा।

%%K t09-tit-7 sub
\VUpk{10.2pt}{40}{पहाड़ पर चढ़ाई।}
\VUsaut{3.0mm}

%%K t09-c1-16-1 p
१६. --- लगभग रोज हम \VUgras{रस्ता देखावे वाला} \textsuperscript{(48)} के
साथे \VUgras{चढ़ाई} के अभ्यास करे जात बानी, आ हमरा बहुते गरब होला जब,
पीठ पर \VUgras{झोरा} \textsuperscript{(47)} आ हाथ में « \VUgras{पहाड़ी
लाठी} », सच्चा \VUgras{सैलानी} \textsuperscript{(46)} नियर हम
\VUgras{चट्टान} \textsuperscript{(45)} पर टूट पड़त बानी।

%%K t09-c1-17-1 p
१७. --- पहाड़ के चोटी से मैदान के रहे वाला लोग चींटी से बड़ ना लागे ;
\VUgras{भेड़ के रेवड़} \textsuperscript{(50)} जे \VUgras{चरागाह}
\textsuperscript{(49)} में चरत बा ऊ बच्चा के खेलौना लागेला।
\VUgras{चीड़} \textsuperscript{(51)} भा \VUgras{काठ के घर}
\textsuperscript{(52)} त हरियरी पर मुसकिल से एगो दाग नियर लागेला।

%%K t09-c1-18-1 p
१८. --- एह घुमाई में हमनी के अकसर \VUgras{पगडंडी} \textsuperscript{(53)}
पर \VUgras{चामोइस के सिकारी} \textsuperscript{(54)} मिलेलें, जे आपन कान्ह
पर ऊ जनावर ढोवत बाड़ें जेकरा के ऊ अभिए मार के आइल बाड़ें।

%%K t09-c1-19-1 p
१९. --- हम आपन पत्ता-संग्रह में \VUgras{फर्न} \textsuperscript{(55)} के
एगो बहुते देखे लायक डाढ़ आ \VUgras{हीथ} \textsuperscript{(57)} के एगो
सुंदर गुलाबी टहनी धर लेनी, जवन हम आपन पिछला घुमाई में एगो छोट
\VUgras{गुफा} \textsuperscript{(56)} के मुहाना पर तोड़ले रहनी।

%%K t09-tit-8 sub
\VUsaut{3.0mm}
\VUcentre{12.0pt}{0}{\textit{दुसरका दृस्य।}}
\VUsaut{3.0mm}

%%K t09-tit-9 sub
\VUpk{11.4pt}{40}{जंगल। --- सिकार।}
\VUsaut{3.0mm}

%%K t09-c2-01-1 p
१. --- कल्ह हम जंगल में एगो बहुते नीक दिन बितवनी। ओहमें रहे वाला
जनावर आ \VUgras{कीड़ा} \textsuperscript{(5)} लोग के बीच के फुरती हमरा
बहुते नीक लागल।

%%K t09-c2-01-2 p
\VUgras{मकड़ी} \textsuperscript{(1)} जे दू गो डाढ़ के बीच आपन
\VUgras{जाला} \textsuperscript{(2)} बुनत बा ताकि गुजरत \VUgras{माछी}
\textsuperscript{(3)} के पकड़ सके, \VUgras{घोंघा} \textsuperscript{(4)} जे
मेहनत से आपन खोल ढोवत बा, सींग वाला भँवरा जे आपन सिकार खोजत बा, मेहनती
चींटी जे \VUgras{चींटी के बमी} \textsuperscript{(6)} लगे बहुते जोस में
बा — ई सब छोट परानी आवत, जात आ काम करत रहल, अजीब लगन से।

%%K t09-c2-02-1 p
२. --- एगो \VUgras{साँप} \textsuperscript{(7)}, जेकरा के हम पहिले
\VUgras{अफई} समझनी, बाकिर जे सादा \VUgras{धामिन} से जादे कुछ ना रहे,
एगो गाछ के ठूँठ के नीचे दुबकल रहे, आ बीच-बीच में आपन पातर छोट कपार
निकालत रहे जवन हरमेसा डँसे खातिर तइयार लागत रहे। हमरा सामने एगो मोटका
\VUgras{बिन-खोल घोंघा} \textsuperscript{(8)} भारी चाल से रेंगत रहे, ओह
सवाद वाला \VUgras{जंगली स्ट्राबेरी} \textsuperscript{(11)} में से एगो
खइला के बाद जे ओहिजा भरपूर उगेला।

%%K t09-c2-03-1 p
३. --- जमीन \VUgras{जंगल के फूल} से बिछल रहे : \VUgras{प्रिमरोज}
\textsuperscript{(9)}, \VUgras{परविंकल} \textsuperscript{(10)}, आ बहुते
दोसर पौधा जेकरा के हम ना जानत रहनी, \VUgras{घास के गुच्छा}
\textsuperscript{(12)} के बीच खिलल रहे।

%%K t09-c2-04-1 p
४. --- ओह इलाका में \VUgras{ट्रफल} लगभग ओतने आम बा जेतना
\VUgras{कुकुरमुत्ता} \textsuperscript{(14)}, आ एगो बनरछक के
\VUgras{सूअर के बच्चा} \textsuperscript{(13)} चटोरपन से खोजत रहे कि ओकरा
कुछ मिल जाव कि ना।

%%K t09-tit-10 sub
\VUsaut{3.0mm}
\VUpk{10.2pt}{40}{चोरी के सिकार।}
\VUsaut{3.0mm}

%%K t09-c2-05-1 p
५. --- हम एगो दृस्य के \VUgras{आखिर} देखनी जवन हमरा बहुते नीक लागल।
आपन \VUgras{बासे कुकुर} \textsuperscript{(15)} के सूँघे के मदद से
\VUgras{बनरछक} \textsuperscript{(16)} अभिए एगो \VUgras{चोर-सिकारी}
\textsuperscript{(22)} के धर लेले रहलें, ठीक ओह बेरा जब ऊ आपन मारल
\VUgras{सिकार} \textsuperscript{(17)} ढोवे के तइयारी करत रहे। धरा जाए के
बदले आपन सिकार छोड़ल नीक समझ के चोर-सिकारी भाग गइल रहे, आ
\VUgras{खरहा} \textsuperscript{(18)}, \VUgras{हिरनी} \textsuperscript{(19)},
\VUgras{छोट हिरन} \textsuperscript{(20)} आ \VUgras{जंगली सूअर}
\textsuperscript{(21)} घास पर पड़ल रह गइल आ बनरछक के खुस कर देलस।

%%K t09-c2-06-1 p
६. --- जंगल में जेतना तरह के गाछ बा ऊ देख के अचरज होला। \VUgras{बीच
के गाछ} \textsuperscript{(23)} आ \VUgras{भोजपत्र} \textsuperscript{(26)},
\VUgras{चीड़} जे \VUgras{गोंद} \textsuperscript{(27)} देला, \VUgras{बलूत}
\textsuperscript{(29)} आ बहुते दोसर गाछ आपन डाढ़ आपस में मिला देले बा।

%%K t09-c2-06-2 p
\VUgras{आइवी} \textsuperscript{(24)}, जे ऊँच गाछ के ठूँठ से चिपकल बा,
\VUgras{जंगल} \textsuperscript{(25)} के चमकत हरियरी से रँग देत बा, जवन ओह
\VUgras{काई} \textsuperscript{(31)} के हल्का पियरइल हरियर रंग से नीक मेल
खात बा जवना में \VUgras{कठफोड़वा} \textsuperscript{(30)} कीड़ा खोजत बा।

%%K t09-tit-11 sub
\VUsaut{3.0mm}
\VUpk{10.2pt}{40}{जंगल के नजारा।}
\VUsaut{3.0mm}

%%K t09-c2-07-1 p
७. --- पुरनका बलूत हमरा मन भा गइल। ओकर मोटका \VUgras{जड़}
\textsuperscript{(32)} ऐंठल बा आ आधा जमीन से बाहर निकलल बा, जवना से ओकरा
में ताकत आ जोर के सुंदर रूप आ जाला, आ हम आपने से पूछत बानी कि
\VUgras{मालिक} \textsuperscript{(35)} एकरा के कटवा के \VUgras{आरी}
\textsuperscript{(34)} के हवाले करे खातिर, जवन \VUgras{लकड़हारा}
\textsuperscript{(33)} के ह, कइसे तइयार हो जालें। बेचारा \VUgras{ठूँठ} \textsuperscript{(36)},
जेकर \VUgras{छाल} \textsuperscript{(37)} छिलल बा भा जेकरा के
\VUgras{टाँगी} \textsuperscript{(38)} चीर देले बा, जवन \VUgras{दिहाड़ी
मजूर} \textsuperscript{(39)} के ह, हमरा दुख देला।

%%K t09-c2-08-1 p
८. --- लगहीं एगो बुढ़उ \VUgras{कोयला बनावे वाला} \textsuperscript{(41)}
आपन \VUgras{बेलचा} \textsuperscript{(42)} से कोयला के \VUgras{ढेर} लगवले
बाड़ें, आ एगो बुढ़िया कटल गाछ के टहनी के \VUgras{सूखल लकड़ी} के
\VUgras{गट्ठर} \textsuperscript{(40)} ले आइल बाड़ी।

%%K t09-tit-12 sub
\VUsaut{3.0mm}
\VUpk{10.2pt}{40}{सिकार।}
\VUsaut{3.0mm}

%%K t09-c2-09-1 p
९. --- हम कुछ पल सुस्तावे खातिर \VUgras{बनरछक के घर} \textsuperscript{(46)}
जाए चाहत रहनी, बाकिर जब हम \VUgras{पोखरा} \textsuperscript{(43)} लगे
पहुँचनी, जवना पर सुंदर \VUgras{राजहंस} \textsuperscript{(45)} तइरत रहे,
हम देखनी कि हाले के \VUgras{बाढ़} \textsuperscript{(44)} रस्ता के सचहूँ
\VUgras{दलदल} बना देले बा। हमरा घूम के जाए के पड़ल, आ \VUgras{देवदार}
\textsuperscript{(49)}, \VUgras{सफेदा} \textsuperscript{(48)} आ
\VUgras{एल्म} \textsuperscript{(47)} लगे से गुजरत, जे जंगल के किनारे बा,
हम देखनी कि \VUgras{झाड़ी} \textsuperscript{(54)} से एगो सुंदर
\VUgras{बारहसिंगा} \textsuperscript{(50)} निकलल, जेकर पाछे जिद्दी
\VUgras{सिकारी कुकुर के झुंड} \textsuperscript{(51)} लागल रहे, आ जेकरा के
\VUgras{बिगुल} \textsuperscript{(53)} भी उकसावत रहे जवन
\VUgras{घोड़सवार सिकारी} \textsuperscript{(52)} के रहे। बेचारा जनावर पूरा
थक गइल रहे। ऊ हमरा से कुछे डेग दूर मरत गिर गइल।

%%K t09-c2-10-1 p
१०. --- बनरछक के बेटा, जेकरा के \VUgras{सिकार के सोर} खींच लइले रहे,
घर से बाहर निकलल आ हमनी के दुनो जंगल में लवट गइनी जा। ऊ एगो पुरनका
बलूत के डाढ़ पर \VUgras{मैगपाई} \textsuperscript{(56)} देखले रहे। ओकरा
लागल कि ओकर घोंसला \VUgras{पत्ता} \textsuperscript{(58)} में छिपल होई, आ
ऊ घोंसला निकाले चाहत रहे।

%%K t09-c2-10-2 p
\VUgras{गिलहरी} \textsuperscript{(61)} नियर फुरतीला होके ऊ एगो
\VUgras{डाढ़} \textsuperscript{(55)} से दोसरा पर चढ़ल, बाकिर ओकरा कुछ ना
मिलल आ ऊ खाली एगो पुरनका \VUgras{कउआ} \textsuperscript{(60)} के उड़ा
पवलस जे एगो \VUgras{टहनी} \textsuperscript{(57)} पर बइठल रहे।
